\section{Alphabete, Wörter, Sprachen}
\begin{mydefinition}{\textbf{Alphabet}}\label{mydefinition:Alphabet}

Eine endliche nichtleere Menge $\Sigma$ heisst \textit{Alphabet}. Die \textit{Elemente eines Alphabets} werden \textit{Buchstaben} (\textit{Zeichen, Symbole}) genannt.
\end{mydefinition}
\cref{mydefinition:Alphabet} entspricht unserer intuitiven Vorstellung eines Alphabets:
\begin{itemize}
	\item Um Text darstellen zu können, muss ein Alphabet mindestens ein Symbol enthalten ($\to$ nichtleer).
	\item Damit man sich auf einen Satz von Zeichen einigen kann, darf das Alphabet nicht unendlich gross sein ($\to$ endlich).
\end{itemize}

Wir listen nun einige der Alphabete auf, die in der Mathematik und Informatik häufig verwendet werden:
\begin{itemize}
	\item $\Sigma_{\text{bool}} = \lrc{0,1}$ ist das Boole'sche Alphabet, mit dem digitale Rechner arbeiten.
	\item $\Sigma_{\text{lat}} = \lrc{a,b,c,...,z,A,B,...,Z}$ ist das lateinische Alphabet.
	\item $\Sigma_{\text{Tastatur}} = \lrc{a,b,c,...,z,A,B,C,...,Z,\text{\textvisiblespace}\,,>,<,(,),...,\#,?,!}$ ist das Alphabet aller Symbole, die mit der englischen Tastatur getippt werden können. Dabei ist \textvisiblespace\ das Symbol für das Leerzeichen / Leersymbol.
	\item $\Sigma_{\text{greek}} = \lrc{\alpha, \beta, \gamma, ..., \omega, A, B, \Gamma, ..., \Omega}$ ist das griechische Alphabet.
	\item $\Sigma_m = \lrc{0,1,2,...,m-1}$ für jedes $m\in\Nunit$ ($\N$ ohne $0$) ist ein Alphabet für die $m$-adische Darstellung von Zahlen.
\end{itemize}
\textit{Wörter} werden wir als endliche Folgen von Buchstaben ansehen.
\begin{mydefinition}{\textbf{Wort}}\label{mydefinition:Wort}
\begin{itemize}
	\item Sei $\Sigma$ ein Alphabet. Ein \textit{Wort} über $\Sigma$ ist eine endliche (möglicherweise leere) Folge von Buchstaben aus $\Sigma$.
	\item Das \textit{leere Wort} $\lambda$ ist die leere Buchstabenfolge.
	\item Die \textit{Länge} $|w|$ eines Wortes $w$ ist die Länge des Wortes als Folge, das heisst die Anzahl der Vorkommen von Buchstaben in $w$.
\end{itemize}
\end{mydefinition}

\begin{myexample}
\begin{aenum}
	\item $w := 1,0,0,1,0$ ist ein Wort über dem Alphabet $\Sigma_{\text{bool}}$ und $|w| = 5$, da $w$ eine Folge von $5$ Buchstaben ist.
	\item $u := M,\text{\textvisiblespace}\,, j$ ist ein Wort über $\Sigma_{\text{Tastatur}}$ und $|u| = 3$, da $u$ eine Folge von $3$ Buchstaben ist.
	\item Das leere Wort $\lambda$ ist ein Wort über jedem Alphabet und es gilt $\abs{\lambda} = 0$.
\end{aenum}
\end{myexample}

\begin{myexercise}
	Was ist $\Sigma_{10}$?
	\begin{myanswer}
		\noindent
		\begin{align*}
			\Sigma_{10} = \lrc{0,1,2,...,9}
		\end{align*}
		ist das Alphabet für die Darstellung von Zahlen im Dezimalsystem.
	\end{myanswer}
\end{myexercise}

\begin{mydefinition}{\textbf{Sigma Stern}}

Sei $\Sigma$ ein Alphabet. $\Sigma^*$ (gesprochen: \textit{Sigma Stern}) ist die Menge aller Wörter über $\Sigma$, also die Menge aller endlichen Folgen von Symbolen aus $\Sigma$.
\end{mydefinition}

\begin{myremark}
	\begin{aenum}
		\item Wir werden im Folgenden Wörter ohne Kommas schreiben. Anstelle von $1,0,0,1,0$ werden wir lediglich $10010$ schreiben. Allgemein, werden wir anstelle von $x_1,x_2,...,x_n$ einfach $x_1x_2...x_n$ schreiben.
		\item In der deutschen Sprache existieren die zwei Ausdrücke \enquote{Wörter} und \enquote{Worte}. Dies sind keine Synonyme. \textit{Wörter} bezeichnet den Plural von \textit{Wort} (\enquote{Im Duden stehen viele Wörter}). Der Ausdruck \textit{Worte} bezieht sich auf Gedankenkonstrukte (\enquote{Sie sprach weise Worte}).
	\end{aenum}
\end{myremark}


\begin{myexercise}\label{myexercise:2.6}
	Wie sieht die Menge $\lrc{1}^*$ aus und wie die Menge $\lr{\Sigma_{\text{bool}}}^*$?
	\begin{myanswer}
		\begin{align*}
			\lrc{1}^* & = \lrc{\lambda, 1, 11, 111, 1111, 11111, ...} =                                         \\
			          & = \lrc{\lambda}\cup\setcm{x_1x_2...x_i}{i\in\N,\; x_j = 1\;\text{für}\; j = 1,2,..., i}
		\end{align*}
		\begin{align*}
			\lr{\Sigma_{\text{bool}}}^* = \lrc{0,1}^* & = \lrc{\lambda, 0, 1,, 00, 01, 10, 11, 000, 001, 010, 100, 011, ...} =                                       \\
			                                          & = \lrc{\lambda}\cup \setcm{x_1x_2...x_i}{i\in\N,\; x_j\in\Sigma_{\text{bool}} \;\text{für}\; j = 1,2,..., i}
		\end{align*}
	\end{myanswer}
\end{myexercise}

\begin{myexercise}
	In \cref{mydefinition:Wort} haben wir ein Wort als endliche Folge von Buchstaben definiert. Nun enthält zum Beispiel die Menge $\lrc{1}^*$ in \cref{myexercise:2.6} aber Wörter beliebiger Länge. Erklären Sie, warum dies kein Widerspruch ist.
\end{myexercise}
\begin{myanswer}
	\noindent
	Anhand dieser Frage lässt sich der Unterschied zwischen den Begriffen \textit{unbeschränkt} und \textit{unendlich} schön verdeutlichen. Betrachten wir dazu exemplarisch nochmals die Menge $\lrc{1}^*$ aus \cref{myexercise:2.6}:
	\begin{align*}
		\lrc{1}^* & = \lrc{\lambda, 1, 11, 111, 1111, 11111, \ldots} =                                              \\
		          & = \lrc{\lambda}\cup \setcm{x_1x_2\ldots x_i}{i\in\N,\; x_j = 1\;\text{für}\; j = 1,2,\ldots, i}
	\end{align*}
	\begin{itemize}
		\item Für jede natürliche Zahl $n\in\N$, existiert in $\lrc{1}^*$ ein Wort $w$ mit $\abs{w} \geq n$, da für jedes $n\in\N$ (insbesondere) auch das Wort $w:=1^n$ in der Menge $\lrc{1}^*$ enthalten ist und $\abs{w}\geq n$. Damit gibt es \textbf{keine obere Schranke} für die Länge der Wörter in $\lrc{1}^*$.
		\item Jedes Wort $w\in\lrc{1}^*$ hat die Form $w = 1^n$ für eine natürliche Zahl $n$. Damit hat jedes Wort in der Menge $\lrc{1}^*$ eine endliche Länge.
	\end{itemize}
	Die Länge der Wörter in $\lrc{1}^*$ können also nicht nach oben beschränkt werden, trotzdem hat jedes Wort in $\lrc{1}^*$ eine endliche Länge und ist somit tatsächlich ein Wort im Sinne von Definition \cref{mydefinition:Wort}.
\end{myanswer}


Mit der folgenden Definition wollen wir den Begriff \textit{Teilwort}, für den wir eine gute Intuition haben, formalisieren. Ein Teilwort eines Wortes $x$ ist ein zusammenhängender Teil von $x$.
\begin{mydefinition}{\textbf{Teilwörter}}\label{mydefinition:Teilwort}

Seien $u,w\in\Sigma^*$ für ein Alphabet $\Sigma$.
\begin{itemize}
	\item $v$ heisst \textbf{Teilwort} von $w$ $\Leftrightarrow$ es existieren $x,y\in\Sigma^*$, so dass $w = xvy$
	\item $v$ heisst \textbf{Präfix} von $w$ $\Leftrightarrow$ es existiert $x\in\Sigma^*$, so dass $w = vx$
	\item $v$ heisst \textbf{Suffix} von $w$ $\Leftrightarrow$ es existiert $x\in\Sigma^*$, so dass $w = xv$
	\item $v\neq \lambda$ heisst \textbf{echtes} Teilwort (Präfix, Suffix) von $w$ $\Leftrightarrow$ $v\neq w$ und $v$ ein Teilwort (Präfix, Suffix) von $w$ ist
\end{itemize}
\end{mydefinition}

\begin{myexample}
Sei $\Sigma = \lrc{a,b,c}$. Das Wort $abc$ ist ein echtes Teilwort, echtes Präfix und ein echtes Suffix von $(abc)^3$. Das leere Wort $\lambda$ ist Teilwort von jedem Wort. Jedes Wort ist Teilwort von sich selbst (aber kein echtes Teilwort).
\end{myexample}

\begin{mychallenge}
	Sei $\Sigma = \lrc{a_1,a_2,...,a_m}$ ein Alphabet der Grösse $m\in\Nunit$.
	\begin{aenum}
		\item Welche Bedingungen muss ein Wort $x$ der Länge $n\leq m$ über $\Sigma$ erfüllen, damit es die maximale Anzahl von verschiedenen Teilwörtern hat?
		\item Wie viele verschiedene Teilwörter hat ein solches $x$?
	\end{aenum}
\end{mychallenge}
\begin{myanswer}
	\begin{aenum}
		\item Die Bedingung ist, dass $x$ (Länge $n$) aus $n$ verschiedenen Buchstaben besteht (also aus lauter verschiedenen Buchstaben). Dann (und nur dann) sind offensichtlich alle möglichen Teilwörter paarweise verschieden und wir erreichen die maximale Anzahl von Teilwörtern.
		Ein solches Wort ist beispielsweise $w = a_1a_2...a_n$.
		\item Es genügt also, die Anzahl der Teilwörter von $w$ zu zählen. Wir werden die Anzahl der Teilwörter der Länge $i = 1, i = 2$ bis zur Länge $i = n$ zählen.\\
		\underline{$i = 1$}:
		\begin{align*}
			1:\quad     & \blue{\underline{a_1}}a_2a_3a_4...a_{n-3}a_{n-2}a_{n-1}a_n \\
			2:\quad     & a_1\blue{\underline{a_2}}a_3a_4...a_{n-3}a_{n-2}a_{n-1}a_n \\
			3:\quad     & a_1a_2\blue{\underline{a_3}}a_4...a_{n-3}a_{n-2}a_{n-1}a_n \\
			            & ...                                                        \\
			(n-3):\quad & a_1a_2a_3a_4...\blue{\underline{a_{n-3}}}a_{n-2}a_{n-1}a_n \\
			(n-2):\quad & a_1a_2a_3a_4...a_{n-3}\blue{\underline{a_{n-2}}}a_{n-1}a_n \\
			(n-1):\quad & a_1a_2a_3a_4...a_{n-3}a_{n-2}\blue{\underline{a_{n-1}}}a_n \\
			n: \quad    & a_1a_2a_3a_4...a_{n-3}a_{n-2}a_{n-1}\blue{\underline{a_n}}
		\end{align*}
		Offensichtlich gibt es $n$ viele Teilwörter der Länge $1$, nämlich genau die Teilwörter $\blue{a_1},\blue{a_2},...,\blue{a_n}$. Wir haben jeweils die Anfangsposition des Teilwortes in $w$ unterstrichen.
		\underline{$i = 2$}:
		\begin{align*}
			1: \quad    & \blue{\underline{a_1}a_2}a_3a_4...a_{n-3}a_{n-2}a_{n-1}a_n \\
			2: \quad    & a_1\blue{\underline{a_2}a_3}a_4...a_{n-3}a_{n-2}a_{n-1}a_n \\
			3: \quad    & a_1a_2\blue{\underline{a_3}a_4}...a_{n-3}a_{n-2}a_{n-1}a_n \\
			            & ...                                                        \\
			(n-3):\quad & a_1a_2a_3a_4...\blue{\underline{a_{n-3}}a_{n-2}}a_{n-1}a_n \\
			(n-2):\quad & a_1a_2a_3a_4...a_{n-3}\blue{\underline{a_{n-2}}a_{n-1}}a_n \\
			(n-1):\quad & a_1a_2a_3a_4...a_{n-3}a_{n-2}\blue{\underline{a_{n-1}}a_n}
		\end{align*}
		Es gibt $n-1$ viele Teilwörter der Länge $2$. Man beachte, dass sich die Anfangsposition der Teilwörter in $w$ um eine Position nach links verschoben hat, da das letzte Teilwort (ganz rechts) der Länge $2$ beim zweitletzten Symbol in $w$ beginnen muss.\\
		\underline{$i = 3$}:
		\begin{align*}
			1: \quad    & \blue{\underline{a_1}a_2}a_3a_4...a_{n-3}a_{n-2}a_{n-1}a_n \\
			2: \quad    & a_1\blue{\underline{a_2}a_3a_4}...a_{n-3}a_{n-2}a_{n-1}a_n \\
			            & ...                                                        \\
			(n-3):\quad & a_1a_2a_3a_4...\blue{\underline{a_{n-3}}a_{n-2}a_{n-1}}a_n \\
			(n-2):\quad & a_1a_2a_3a_4...a_{n-3}\blue{\underline{a_{n-2}}a_{n-1}a_n}
		\end{align*}
		Es gibt $n-2$ viele Teilwörter der Länge $3$.\\

		\noindent
		Es ist nun nicht mehr schwierig zu sehen, dass es für ein Teilwort der Länge $i\in\lrc{1,2,...,n}$ genau $n-i+1$ mögliche Anfangspositionen in $w$ gibt, da die Positionen $n-i+2,3,...,n$ nicht möglich sind (\enquote{zu wenig Platz}). Nun summieren wir die Anzahl aller unterschiedlichen (nicht leeren) Teilwörter von $w$ der Längen $i=1,2,...n$ auf:
		\begin{align*}
			\sum_{i=1}^n (n-i+1) = n(n+1)-\sum_{i=1}^n i = n(n+1) - \frac{n(n+1)}{2}=\frac{n(n+1)}{2}
		\end{align*}
		Hinzu kommt noch das leere Wort (welches Teilwort jedes Wortes ist) und als Resultat finden wir, dass es
		\begin{align*}
			\frac{n(n+1)}{2} + 1
		\end{align*}
		viele verschiedene Teilwörter von $w$ gibt.
	\end{aenum}
\end{myanswer}

